\chapter{Sequences, cycles and external change}\label{ch:cycles}
The simplest accounting theorem in EBU is also one of the most important.
If every action begins where the previous one ended, the intermediate
potentials cancel. A surprisingly large number of incorrect issuance rules
fail because they stop using those live intermediate states.

\section{Undriven telescoping}
Let $x^0,x^1,\ldots,x^N$ be a declared sequence under one fixed potential,
with action values $E_k=V(x^k)-V(x^{k+1})-C_k$. Then
\begin{equation}\label{eq:telescope}
 \sum_{k=0}^{N-1}E_k
 =V(x^0)-V(x^N)-\sum_{k=0}^{N-1}C_k.
\end{equation}
To prove it, write the first two differences: $V(x^0)-V(x^1)$ and
$V(x^1)-V(x^2)$. Their shared term cancels. Repeating that cancellation
leaves only the first and last endpoints. Neither convexity nor Gaussian
geometry is required for this algebraic identity.

For a closed physical cycle $x^N=x^0$, the total is $-\sum C_k$.
If all burdens are nonnegative, it cannot be positive. With $C=0$ it is
zero, not necessarily negative. A reversible movement can return every
represented coordinate to its start without making the account grow.

\section{A live-state calculation}
Use $(18,2)$ with reference $(10,10)$ and scales $(2,2)$. Transfer two
units, then three, then one. The successive potentials are
\begin{equation}
 16,\quad9,\quad9/4,\quad1.
\end{equation}
The action values are $7$, $27/4$ and $5/4$, whose sum is fifteen.
That equals $16-1$. Quoting all three against the first state instead
would replace the last two live values with unrelated counterfactuals.

Splitting a transfer preserves the field total when the resulting sequence
has the same endpoint. It need not preserve separate activation costs,
physical delays or intervening disturbances. A theorem about endpoint
differences is not a theorem that every way of organizing work costs the same.

\section{Forcing before the action epoch}
The new programme proposes an external change from $x_t$ to $z_t$ before
actors evaluate their candidates. Define
\begin{equation}
 D_{\mathrm{ext},t}=V(z_t)-V(x_t),\qquad
 E_t=V(z_t)-V(x_{t+1})-C_t.
\end{equation}
Add and subtract $V(x_t)$ in the second expression and sum:
\begin{equation}\label{eq:driven}
 \sum_{t=0}^{N-1}E_t
 =V(x_0)-V(x_N)+\sum_{t=0}^{N-1}D_{\mathrm{ext},t}
                  -\sum_{t=0}^{N-1}C_t.
\end{equation}
An actor who does nothing has zero field value at the shared post-forcing
baseline. Nature's own change appears in $D_{\mathrm{ext}}$, not in an
idle actor's income. The external term can be negative: forcing may move
the system closer to the reference.

The old D0 foundation uses a frozen-drive no-action successor. Its immediate
cancellation reasoning is related, but its event order is not identical to
this proposed forcing-first profile. An implementation needs an explicit
profile contract rather than renaming one variable and claiming equivalence.

\section{A forcing example with exact signs}
Begin at reference $(10,10)$, where $V=0$. A conservative external shift
produces $(14,6)$, where $V=4$. An actor then transfers two units back,
ending at $(12,8)$ with $V=1$. The external audit is $+4$ and the actor's
field value is $+3$. The total physical potential change is only $+1$;
confusing it with actor value would get both the baseline and the sign wrong.

If a later external shift takes $(12,8)$ to $(10,10)$, its audit is $-1$.
No actor performed that improvement. Crediting it to the next idle actor
would violate the declared separation even if a global sum happened to close.

\section{Changing the potential is another accounted event}
If parameters vary, simple telescoping under one $V$ no longer applies.
Writing $V_t$ for the version at event $t$, parameter changes create terms
such as $V_{t+1}(x)-V_t(x)$ even at unchanged physical state. A recalibration
can therefore create apparent value unless its effect is separated from
physical actions. The first model freezes references and scales precisely
to avoid silently mixing those effects.

The historical no-overissue theorem retains its original one-action and
audit assumptions. This chapter derives a general endpoint identity; it
does not retroactively broaden that source theorem's settlement scope.
For simultaneous groups, the joint transition must be valued once and any
division needs the distinct treatment of Chapters~\ref{ch:groups} and~\ref{ch:attribution}.

\section*{Chapter recap}
Live endpoints, fixed parameters and complete external accounting make the
sum reconstructible. Closed-cycle closure rules out one accounting exploit.
It does not prove that actors prefer recovery over cycling or refusal.
